\documentclass[aspectratio=169,11pt,t]{beamer}
\input{../../../theme/ksa-cs-v2.tex}
\usepackage{fontawesome5}
\newcommand{\paperlink}[1]{\href{#1}{\textcolor{ksablue}{\faLink}}}
% GLM 차시 (2026-11-27): 앤드류 응 CS229 강의 노트 Part III (Generalized Linear Models) 을 덱 스타일로.
% 숫자는 data.py (assert), 그림은 make_figs.py.
\setbeamerfont{author}{size=\normalsize}
\setbeamerfont{institute}{size=\small}
\graphicspath{{figs/}{../../../assets/}}
\usetikzlibrary{positioning,arrows.meta,calc}

\renewcommand{\ksacredit}{Machine Learning and Deep Learning \,\textperiodcentered\, GLM}
\title{일반화 선형 모델 (GLM)}
\subtitle{포아송 회귀, 그리고 모델들을 하나로 묶는 식}
\author{Machine Learning and Deep Learning}
\institute{KSA 2027-Spring}
\date{Sumin Han (suminhan@ksa.hs.kr)}
\begin{document}
{\setbeamertemplate{footline}{}\begin{frame}[c,noframenumbering]\titlepage\end{frame}}

\section{질문}
% --- 오늘의 질문: 개수일 때
\begin{frame}{오늘의 질문: $y$ 가 개수일 때는?}
\vfill
\centering
{\large 지금까지: \;실수 $\to$ 선형회귀 \;·\; 0 / 1 $\to$ 로지스틱회귀 \;·\; 여러 클래스 $\to$ 소프트맥스}\\[1.4em]
{\LARGE 이 팀은 이번 경기에서 \kq{몇 골} 넣을까? \;(0, 1, 2, \dots)}\\[1.2em]
{\large 직선 $z = w_0 + w_1 x$ 에 무엇을 씌우고, 손실은 무엇으로 할까?}\\[0.5em]
{\large 그리고: 이 모델들을 \kb{하나로 묶는 식}이 있을까?}
\vfill
\note{지금까지 세 가지 모델을 배웠습니다. 점수처럼 아무 실수인 y는 선형회귀, 체스 승패처럼 0 아니면 1인 y는 로지스틱회귀, 붓꽃처럼 여러 클래스 중 하나인 y는 소프트맥스 회귀로 맞혔습니다. 오늘은 하나를 더 해 봅니다. 개수입니다. 월드컵에서 이 팀이 이번 경기에 몇 골을 넣을까요? 0, 1, 2로 세는 값입니다. 1주차부터 써 온 직선 z에 무엇을 씌우고, 손실은 무엇으로 해야 할까요? 그리고 마지막에 이 모델들을 하나로 묶는 식이 있는지 봅니다.}
\end{frame}

\section{개수}
% --- 포아송 분포
\begin{frame}{개수의 분포: 포아송}
\vfill
\centering
{\Large\[ P(y) = \frac{\lambda^{y}\,e^{-\lambda}}{y!} \qquad y = 0, 1, 2, \dots \qquad \text{평균} = \lambda \]}\vskip-0.2em
\includegraphics[width=0.43\linewidth]{glm_pois_pmf.png}\par\vskip0em
{\large 평균 $\lambda$ 하나가 모양을 정한다 \;·\; $\lambda$ 가 클수록 \kb{오른쪽으로, 넓게}}\\[0em]
{\small 데이터: 월드컵 본선 1068경기, 팀마다 한 줄 (2136줄) \quad $x$ = Elo 차이 (내 팀 $-$ 상대)}
\vfill
\note{데이터는 1930년 첫 대회부터 올해 2026년 대회까지 월드컵 본선 23개 대회 1068경기입니다. 팀마다 한 경기에 넣은 골을 한 줄로 적어 2136줄입니다. x는 두 팀의 실력 차이, Elo 점수 차이입니다. Elo는 체스 레이팅과 같은 방식으로, 국가대표 경기 결과로 매긴 점수입니다. 골 수는 0, 1, 2로 세는 개수입니다. 개수에 쓰는 분포가 포아송입니다. 평균이 람다일 때 y개가 나올 확률은 람다의 y제곱 곱하기 e의 마이너스 람다 나누기 y 팩토리얼입니다. 평균 람다 하나만 정하면 모양이 다 정해집니다. 그림은 뒤에서 학습한 모델이 예측한 세 경우입니다. 빨강은 Elo가 300점 뒤진 약팀, 평균 0.74골로 0골일 확률이 거의 절반입니다. 회색은 실력이 비슷한 팀, 평균 1.34골입니다. 파랑은 300점 앞선 강팀, 평균 2.40골로 2골, 3골도 흔합니다. 람다가 커질수록 오른쪽으로 가면서 넓게 퍼집니다. 그러니 x로 람다만 잘 맞히면 됩니다.}
\end{frame}

% --- 포아송 회귀 모델
\begin{frame}{포아송 회귀: 평균 $\lambda$ 를 직선으로}
\vfill
\centering
{\large $\lambda = w_0 + w_1 x$ 로 두면? \;$\to$\; 약팀 쪽에서 \kq{음수 골} (평균 골은 0 이상이어야)}\\[0.8em]
{\LARGE\[ \hat y = \lambda = e^{\,z}, \qquad z = w_0 + w_1 x \]}\vskip0.4em
{\large $e^{z}$ 는 \kb{언제나 양수} \;·\; $x$ 가 1 늘면 $z$ 는 $w_1$ 늘고 $\hat y$ 는 \kb{$e^{w_1}$ 배}}\\[0.6em]
{\large 로지스틱과 같은 구조: 직선 $z$ 에 함수 하나 ($\sigma$ 대신 $e^{z}$)}\\[0.3em]
{\small 거꾸로 쓰면 $z = \log\lambda$ \;(평균의 로그가 직선) \quad 로지스틱은 $z = \log\frac{p}{1-p}$ (오즈의 로그가 직선)}
\vfill
\note{포아송의 평균 람다를 x로 맞히는 모델을 만들어 봅시다. 1주차처럼 람다를 바로 직선으로 두면, 약팀 쪽에서 직선이 0 아래로 내려가 음수 골을 예측합니다. 평균 골은 0 이상이어야 하죠. 그래서 직선 z에 e를 씌웁니다. 와이햇은 람다, 람다는 e의 z입니다. e의 z는 z가 아무리 작아도 언제나 양수입니다. 그리고 x가 1 늘면 z가 w1 늘고, 와이햇은 e의 w1 배가 됩니다. 더해 가는 게 아니라 곱해 갑니다. 로지스틱회귀와 똑같은 구조입니다. 직선 z에 함수 하나를 씌우는데, 시그모이드 대신 e의 z입니다. 거꾸로 쓰면 z는 log 람다, 평균의 로그가 직선입니다. 로지스틱회귀에서 오즈의 로그가 직선이었던 것과 같은 읽는 법입니다.}
\end{frame}

% --- 포아송 손실 · 기울기
\begin{frame}{포아송 회귀의 손실과 기울기}
\vfill
\centering
{\large\begin{align*}
\text{loss} &= -\log P(y) = -\log\frac{\hat y^{\,y} e^{-\hat y}}{y!} = \hat y - y\log\hat y + \log y! && \text{로지스틱처럼 } -\log P\\[0.3em]
\frac{\partial\,\text{loss}}{\partial z} &= \frac{\partial}{\partial z}\big(e^{z} - y\,z\big) = e^{z} - y = \textcolor{ksared}{\hat y - y} && \hat y = e^{z},\; \log\hat y = z\\[0.3em]
\frac{\partial\,\text{loss}}{\partial w_1} &= (\hat y - y)\,x && z = w_0 + w_1 x
\end{align*}}\vskip0.4em
{\large 예: 비슷한 팀 ($\hat y = 1.34$) 이 3골 \;$\to$\; $\hat y - y = -1.66$ \;$\to$\; $w$ 를 바꿔 예측을 \kb{올린다}}\\[0.5em]
{\large 기울기가 \kq{로지스틱회귀와 똑같은} $(\hat y - y)\,x$ \;$\to$\; 경사하강법 코드에서 \texttt{sigmoid} 를 \texttt{exp} 로만}
\vfill
\note{손실은 로지스틱회귀 때처럼 마이너스 log 확률입니다. 실제로 나온 골 y의 확률이 클수록 손실이 작습니다. 포아송 확률에 마이너스 log를 씌우면 와이햇 빼기 y log 와이햇 더하기 log y 팩토리얼입니다. 마지막 항은 w와 상관없어서 미분하면 사라집니다. 와이햇은 e의 z, log 와이햇은 z이니, 손실은 e의 z 빼기 y z이고, z로 미분하면 e의 z 빼기 y, 곧 와이햇 빼기 y입니다. 연쇄 법칙으로 x를 곱하면 w1으로 미분한 값은 와이햇 빼기 y 곱하기 x. 예를 들어 실력이 비슷해서 1.34골을 예측했는데 실제로 3골을 넣었다면, 와이햇 빼기 y가 마이너스 1.66이라 예측을 올리는 쪽으로 w가 바뀝니다. 놀라운 점은 이 기울기가 로지스틱회귀와 똑같다는 것입니다. 그래서 로지스틱회귀 경사하강법 코드에서 시그모이드를 exp로 바꾸기만 하면 포아송 회귀가 됩니다. 이 걸음을 실제로 되풀이해 봅시다.}
\end{frame}

% --- 포아송 학습 과정
\begin{frame}{학습: 경사하강법을 되풀이하면}
\vfill
\centering
\renewcommand{\arraystretch}{1.45}
{\large\begin{tabular}{l|l}
\kb{시작} & $w_0 = \log 1.42 = 0.35,\; w_1 = 0$ \;$\to$\; 모든 팀을 평균 1.42골로 예측\\
\kb{한 걸음} & 2136줄마다 $\hat y = e^{\,w_0 + w_1 x}$ \;$\to$\; 오차 $\hat y - y$ \;$\to$\; 평균 낸 기울기\\
 & $w_0 \leftarrow w_0 - \alpha\cdot\text{평균}(\hat y - y), \quad w_1 \leftarrow w_1 - \alpha\cdot\text{평균}\big((\hat y - y)\,x\big)$\\
\kb{첫 걸음 값} & $\text{평균}\big((\hat y - y)\,x\big) = -0.87$ \;(Elo 100 단위) \;$\to$\; $w_1$ 이 \kb{커진다}\\
 & 강팀 ($x > 0$) 은 실제로 더 넣어서 $\hat y - y < 0$ \;$\to$\; 기울기가 음수\\
\kb{끝} & 20000번 되풀이 \;$\to$\; $w_0 = 0.289,\; w_1 = 0.195$ (100점당) \;(기울기 $\approx 0$)
\end{tabular}}\\[0.8em]
{\large 로지스틱회귀 코드와 \kb{한 군데만} 다르다: \;\texttt{sigmoid(z)} $\to$ \texttt{exp(z)}}
\vfill
\note{그럼 포아송 회귀는 실제로 어떻게 학습될까요? 로지스틱회귀와 똑같이 경사하강법입니다. 시작은 w0를 평균 골 1.42의 로그, 0.35로, w1을 0으로 둡니다. 그러면 Elo 차이와 상관없이 모든 팀을 1.42골로 예측합니다. 한 걸음은 이렇습니다. 2136줄 각각에서 와이햇, e의 w0 더하기 w1 x를 계산하고, 실제 골을 빼서 오차 와이햇 빼기 y를 구합니다. w0는 오차의 평균만큼, w1은 오차 곱하기 x의 평균만큼 학습률을 곱해 빼 줍니다. 첫 걸음에서 오차 곱하기 x의 평균은 Elo 100점 단위로 마이너스 0.87입니다. 강팀, x가 양수인 팀은 실제로 1.42골보다 많이 넣어서 오차가 음수이고, 약팀은 x도 음수, 오차는 양수라서 곱이 역시 음수입니다. 기울기가 음수이니 w1을 키웁니다. 강할수록 더 넣는다는 쪽으로 곡선이 기울어지는 겁니다. 이 걸음을 2만 번 되풀이하면 기울기가 거의 0이 되고, w0 0.289, w1은 Elo 100점당 0.195, 1점당 0.00195에서 멈춥니다. 다음 쪽 곡선이 바로 이 값입니다. 코드로 보면 로지스틱회귀의 sigmoid를 exp로 바꾼 것뿐입니다.}
\end{frame}

% --- 포아송 회귀 결과
\begin{frame}{월드컵 2136줄로 학습하면}
\vfill
\centering\includegraphics[width=0.58\linewidth]{glm_goals_pois.png}\par\vskip0.2em
{\large 골 $= e^{\,0.289 + 0.00195 \times \text{Elo 차이}}$ \quad $-300$: 0.74골 \;·\; 0: 1.34골 \;·\; $+300$: 2.40골}\\[0.2em]
{\large Elo 100 앞설 때마다 \kb{$\times 1.22$} ($e^{0.195}$) \quad 가장 약한 쪽 ($-575$) 도 \kb{0.43골} (직선은 $-0.17$)}
\vfill
\note{점선은 1주차처럼 그은 직선입니다. 약팀 쪽으로 가면 0 아래로 내려가 음수 골을 예측합니다. 2136줄로 포아송 회귀를 학습하면 넣은 골은 e의 0.289 더하기 0.00195 곱하기 Elo 차이입니다. 실력이 같으면 1.34골, 300점 앞서면 2.40골, 300점 뒤지면 0.74골을 예측합니다. Elo가 100점 앞설 때마다 e의 0.195, 1.22배씩 늘어납니다. 그래서 더해 가는 직선이 아니라 곱해 가는 지수 곡선이고, 검은 점, 실제 평균을 직선보다 잘 따라갑니다. 직선이 음수를 예측했던 가장 약한 쪽, 575점 뒤진 경우도 0.43골로 양수입니다. 학습은 로지스틱회귀 코드에서 시그모이드를 exp로 바꾼 경사하강법으로 했고, sklearn의 PoissonRegressor와 같은 답이 나옵니다.}
\end{frame}

\section{일반화}
% --- 맞힐 값의 종류
\begin{frame}{맞힐 값 $y$ 의 종류가 다르다}
\vfill
\centering\includegraphics[width=\linewidth]{glm_types.png}\par\vskip0.3em
{\large 실수 \;·\; 0 / 1 \;·\; \kb{개수} (0 이상 정수) \;·\; \kb{셋 중 하나} \quad $\to$ 예측도 그 모양이어야 한다}\\[0.2em]
{\small 붓꽃 = Fisher (1936) 150송이}
\vfill
\note{포아송 회귀까지 배웠으니 지금까지의 데이터 넷을 나란히 놓고 비교해 봅시다. 맞히려는 값 y는 종류가 여러 가지였습니다. 1주차의 시험 점수는 아무 실수나 될 수 있습니다. 체스 승패는 0 아니면 1입니다. 오늘 본 월드컵 골 수는 0 이상의 정수, 개수입니다. 붓꽃은 꽃잎 길이로 세 종 중 하나를 맞히는 문제입니다. 예측도 y의 모양에 맞아야 합니다. 확률은 0과 1 사이, 개수는 0 이상, 세 종의 확률은 더해서 1이어야 하죠.}
\end{frame}

% --- 분포 네 개
\begin{frame}{데이터마다 맞는 분포가 있다}
\vfill
\centering
\includegraphics[width=0.95\linewidth]{glm_dists}\\[0.1em]
\renewcommand{\arraystretch}{1.2}
{\large\begin{tabular}{cccc}
\makebox[0.21\linewidth]{$\frac{1}{\sqrt{2\pi}\sigma}\,e^{-(y-\mu)^2/2\sigma^2}$} &
\makebox[0.21\linewidth]{$\phi^{\,y}(1-\phi)^{1-y}$} &
\makebox[0.21\linewidth]{$\lambda^{y}e^{-\lambda}/\,y!$} &
\makebox[0.21\linewidth]{$p(y = k) = \phi_k$}\\
$\sigma = 4.4$ 점 & $\phi = 363/500 = 0.73$ & $\lambda = 1.30$ 골 & $\phi = (0,\ 0.38,\ 0.62)$
\end{tabular}}\\[0.4em]
{\small 막대 $=$ 실제 데이터, 검은 선 $=$ 같은 평균의 분포 }
\vfill
\note{y마다 어떤 모양으로 흩어지는지, 분포를 맞춰 봅시다. 막대가 실제 데이터, 검은 선이 같은 평균을 가진 분포입니다. 점수는 직선 예측에서 벗어난 정도를 30명 모아 보면 종 모양, 정규분포입니다. 표준편차는 4.4점입니다. 체스는 레이팅이 200점 앞선 근처 500판을 모으면 백이 363판 이겼습니다. 1일 확률 파이가 0.73인 베르누이입니다. 월드컵은 실력이 비슷한 경기, Elo 차이가 50점 안인 462줄의 골 수입니다. 평균 1.30골인 포아송과 거의 겹칩니다. 붓꽃은 꽃잎 길이 4.8에서 5.0센티미터인 13송이를 모으면 setosa 0, versicolor 5, virginica 8송이, 범주형 분포입니다. 정규분포는 선형회귀, 베르누이는 로지스틱회귀, 포아송은 포아송 회귀, 범주형은 소프트맥스 회귀가 맡습니다. 분포 하나에 모델 하나씩입니다. 그런데 이 네 모델에는 묘한 공통점이 있습니다.}
\end{frame}

% --- 우연일까
\begin{frame}{네 모델, 기울기가 모두 같다: 우연일까?}
\vfill
\centering
\renewcommand{\arraystretch}{1.5}
{\large\begin{tabular}{l|c|c|c}
모델 & 예측 $\hat y$ & 손실 & $\partial\,\text{loss} / \partial z$\\ \hline
선형회귀 & $z$ & $\tfrac12 (y - \hat y)^2$ & \kq{$\hat y - y$}\\
로지스틱회귀 & $\sigma(z)$ & $-\big[y\log\hat y + (1-y)\log(1-\hat y)\big]$ & \kq{$\hat y - y$}\\
포아송 회귀 & $e^{z}$ & $\hat y - y\log\hat y$ & \kq{$\hat y - y$}\\
소프트맥스 회귀 & $\mathrm{softmax}(z)$ & $-\log\hat y_{\text{정답}}$ & \kq{$\hat y - y$}
\end{tabular}}\\[0.8em]
{\large 함수도 손실도 다 다른데 기울기는 \kb{모두} $\hat y - y$ \;$\to$\; 네 모델을 묶는 \kb{일반화된 식}이 있다}\\[0.3em]
{\small GLM (일반화 선형 모델): \kb{직선} (Linear) 은 그대로, \kb{정규분포가 아닌} $y$ 로 일반화 (Generalized)}
\vfill
\note{오늘 배운 포아송 회귀와 지금까지 배운 세 모델을 나란히 놓아 봅시다. 예측 함수는 z 그대로, 시그모이드, e의 z, 소프트맥스로 다 다릅니다. 손실도 오차 제곱, 교차 엔트로피, 포아송 손실, 다중 교차 엔트로피로 다 다릅니다. 그런데 손실을 z로 미분하면 넷 다 와이햇 빼기 y입니다. 우연일까요? 아닙니다. 이 네 모델은 한 가지 식에서 나옵니다. 그 일반화된 식이 일반화 선형 모델, GLM입니다. 이름이 뜻을 말해 줍니다. 선형회귀는 y가 직선 근처에 정규분포로 흩어져 있다고 봅니다. 직선, Linear는 그대로 두고, 정규분포가 아닌 y까지 다룰 수 있게 넓힌 것이 일반화, Generalized입니다.}
\end{frame}

% --- 지수족
\begin{frame}{네 분포가 같은 꼴로 써진다: 지수족}
\vfill
\centering
{\Large\[ p(y;z) = b(y)\,\exp\!\big(\,\underbrace{z}_{\text{natural parameter}}\, \underbrace{T(y)}_{\text{sufficient statistic}} \;-\; \underbrace{a(z)}_{\text{log partition}}\,\big) \]}\vskip0.1em
\renewcommand{\arraystretch}{1.12}
{\large\begin{tabular}{c|l}
$p(y;z)$ & 손잡이를 $z$ 로 맞춘 분포에서 \kb{$y$ 가 나올 확률}\\
$y$ & 실제로 나온 값 (골 수, 승패, 점수)\\
$z$ & 분포를 정하는 \kb{손잡이} (자연 매개변수)\\
$T(y)$ & $y$ 에서 뽑아 쓰는 값 (오늘은 $y$ 그대로)\\
$b(y)$ & $y$ 만 들어 있는 부분\\
$a(z)$ & $z$ 만 들어 있는 부분 $\to$ 확률의 합을 1 로 맞춘다
\end{tabular}}\\[0.3em]
{\large 정규 · 베르누이 · 포아송 · 범주형이 모두 이 꼴 $=$ \kb{지수족}}
\vfill
\note{네 모델의 분포, 정규, 베르누이, 포아송, 범주형은 정리하면 모두 이 꼴로 쓸 수 있습니다. b of y 곱하기 exp z 곱하기 T of y 빼기 a of z. 기호를 하나씩 봅시다. p of y 세미콜론 z는 손잡이를 z로 맞춘 분포에서 y가 나올 확률입니다. 세미콜론 뒤는 설정값이라는 뜻입니다. y는 실제로 나온 값, 골 수나 승패나 점수입니다. z는 분포를 정하는 손잡이로, 돌리면 평균이 바뀝니다. b of y는 y만, a of z는 z만 들어 있는 부분입니다. 이 꼴로 쓸 수 있는 분포들을 지수족이라고 부릅니다. z를 자연 매개변수, T of y를 충분 통계량, a of z를 로그 분배 함수라고 부릅니다. 이름은 어렵지만 역할은 간단합니다. T, a, b 세 가지를 정하면 분포의 가족 하나가 정해지고, 그 안에서 z를 바꾸면 평균이 다른 분포들이 나옵니다. 예를 들어 베르누이 가족에서 z를 바꾸면 앞면 확률이 다른 동전들이 나옵니다. e의 마이너스 a of z는 확률을 다 더했을 때 1이 되도록 맞추는 몫입니다. 오늘 보는 분포는 범주형만 빼고 T of y가 y 그대로입니다. 이제 실제 분포를 이 꼴에 넣어 봅시다.}
\end{frame}

% --- 지수족
\begin{frame}{분포들을 이 꼴에 넣으면}
\vfill
\centering
{\Large\[ p(y) = \textcolor{gray}{b(y)}\,\cdot\,\exp\!\big(\,\textcolor{ksared}{z}\cdot y \;-\; \textcolor{ksablue}{a(z)}\,\big) \]}\vskip0.3em
\renewcommand{\arraystretch}{1.4}
{\large\begin{tabular}{l|c|c|c|c}
분포 & 원래 식 & \textcolor{gray}{$b(y)$}: $y$ 만 & \textcolor{ksared}{$z$}: $y$ 에 곱함 & \textcolor{ksablue}{$a(z)$}: 빼는 몫\\ \hline
포아송 & $\dfrac{\lambda^{y} e^{-\lambda}}{y!}$ & $\dfrac{1}{y!}$ & \textcolor{ksared}{$\log\lambda$} & \textcolor{ksablue}{$\lambda$}\\
베르누이 & $\phi^{\,y}(1-\phi)^{1-y}$ & $1$ & \textcolor{ksared}{$\log\frac{\phi}{1-\phi}$} & \textcolor{ksablue}{$-\log(1-\phi)$}\\
정규 ($\sigma{=}1$) & $\frac{1}{\sqrt{2\pi}}e^{-(y-\mu)^2/2}$ & $\frac{1}{\sqrt{2\pi}}e^{-y^2/2}$ & \textcolor{ksared}{$\mu$} & \textcolor{ksablue}{$\tfrac12\mu^2$}
\end{tabular}}\\[0.5em]
{\large 포아송 확인: $\lambda^{y} = e^{\,y\log\lambda}$ 이라서 $\frac{\lambda^{y} e^{-\lambda}}{y!} = \frac{1}{y!}\exp(\textcolor{ksared}{\log\lambda}\cdot y - \textcolor{ksablue}{\lambda})$ \quad 이 꼴 $=$ \kb{지수족}}\\[0.2em]
{\small $T(y) = y$ 라서 $z^\top T(y) = z \cdot y$ \;·\; 범주형은 $z$ 가 벡터 (뒤에서)}
\vfill
\note{앞에서 본 분포 식들은 모양이 제각각입니다. 그런데 정리하면 모두 방금 본 꼴로 쓸 수 있습니다. T of y는 y 그대로라서 z 곱하기 y로 적었습니다. 회색 b of y는 y만 들어 있는 부분, 빨간 z는 y에 곱해지는 부분, 파란 a of z는 빼 주는 부분입니다. 가장 쉬운 포아송으로 확인해 봅시다. 람다의 y제곱은 e의 y log 람다와 같습니다. 그러면 람다의 y제곱 곱하기 e의 마이너스 람다 나누기 y 팩토리얼은, 1 나누기 y 팩토리얼 곱하기 exp log 람다 곱하기 y 빼기 람다입니다. b of y는 1 나누기 y 팩토리얼, z는 log 람다, a는 람다입니다. 베르누이는 y log 파이 더하기 1 빼기 y log 1 빼기 파이를 y로 묶으면 z가 로그 오즈, 정규분포는 제곱을 전개하면 z가 뮤 그대로입니다. 앞 쪽의 붓꽃, 범주형 분포도 지수족인데 z가 숫자 하나가 아니라 클래스마다 하나씩인 벡터라서 표에서는 뺐고, 뒤의 소프트맥스에서 다시 나옵니다. 포아송이면 z가 평균 람다를 정하고, 베르누이면 앞면 확률 파이를 정합니다.}
\end{frame}

% --- GLM 세 가정
\begin{frame}{GLM: 세 가지만 정하면 된다}
\vfill
\centering
\renewcommand{\arraystretch}{1.6}
{\Large\begin{tabular}{rl}
\textbf{1.} & $y \mid x \;\sim\;$ \kb{지수족}$(z)$ \quad {\large ($y$ 의 분포를 고른다)}\\
\textbf{2.} & 예측 $\hat y = E[\,y \mid x\,]$ \quad {\large (그 분포의 \kb{평균})}\\
\textbf{3.} & $z = w_0 + w_1 x$ \quad {\large (매개변수 $=$ \kb{직선}, 가정이라기보다 \kq{설계 선택})}
\end{tabular}}\\[0.8em]
{\large 2번 확인 (로지스틱): \;$E[\,y \mid x\,] = 0 \cdot P(y{=}0) + 1 \cdot P(y{=}1) = P(y{=}1 \mid x) = \hat y$ \;\checkmark}\\[0.3em]
{\small GLM = Generalized Linear Model (일반화 선형 모델) \quad $x$ 가 여럿이면 $z = w \cdot x + b$}
\vfill
\note{GLM은 모델을 만드는 요리법이고, 세 가지를 정합니다. 첫째, x가 주어졌을 때 y가 지수족의 어떤 분포를 따른다고 고릅니다. 자연 매개변수는 z입니다. 둘째, 예측 와이햇은 그 분포의 평균, x가 주어졌을 때 y의 기댓값으로 합니다. 로지스틱회귀로 확인해 봅시다. y가 0 아니면 1이니 기댓값은 0 곱하기 0일 확률 더하기 1 곱하기 1일 확률, 곧 1일 확률입니다. 로지스틱회귀의 와이햇이 바로 이 확률이었으니 조건을 맞춥니다. 셋째, 자연 매개변수 z를 x의 직선으로 둡니다. 강의 노트는 이것이 가정이라기보다 설계 선택이라고 말합니다. 왜 꼭 직선이어야 하는지 근거가 있는 것은 아니지만, 이렇게 두면 학습이 아주 쉬워집니다. 이 세 가지로 지금까지의 모델이 어떻게 나오는지 봅시다.}
\end{frame}

% --- 유도: 선형회귀 · 로지스틱회귀
\begin{frame}{GLM 으로 유도하면: 선형회귀와 로지스틱회귀 (포아송도 같다)}
\vskip-0.3em
\begin{columns}[T]
\begin{column}{0.48\textwidth}
\centering{\large\kb{$y$ 가 실수: 정규}}
{\large\begin{align*}
\hat y &= E[\,y \mid x\,] && \text{2번}\\
&= \mu && \text{정규의 평균}\\
&= z && z = \mu\\
&= w_0 + w_1 x && \text{3번}
\end{align*}}
\centering{\large $=$ \kb{선형회귀}}
\end{column}
\begin{column}{0.5\textwidth}
\centering{\large\kb{$y$ 가 0 / 1: 베르누이}}
{\large\begin{align*}
\hat y &= E[\,y \mid x\,] && \text{2번}\\
&= \phi && \text{베르누이의 평균}\\
&= \dfrac{1}{1 + e^{-z}} && z = \text{로그 오즈를 풂}\\
&= \dfrac{1}{1 + e^{-(w_0 + w_1 x)}} && \text{3번}
\end{align*}}
\centering{\large $=$ \kb{로지스틱회귀}}
\end{column}
\end{columns}\vskip0.1em
\centering
{\large 시그모이드는 S 자라서 \textbf{고른} 것이 아니라, $y$ 가 0 / 1 이라서 \kb{나온} 것}\\[0.3em]
{\small 체스: 백이 200점 높으면 $z = 0.09 + 0.0047 \times 200 = 1.03$ \;$\to$\; $\sigma(1.03) = 74\%$}\\[0.1em]
{\small $z$ 를 평균으로 바꾸는 $g$ $=$ canonical response function (정규: 그대로, 베르누이: $\sigma$)}
\vfill
\note{세 가지 재료로 선형회귀를 유도해 봅시다. y가 실수라서 정규분포를 고릅니다. 와이햇은 2번에 따라 y의 기댓값, 정규분포의 기댓값은 뮤, 지수족 꼴에서 뮤는 z 그대로였고, 3번에 따라 z는 w0 더하기 w1 x입니다. 와이햇이 직선, 선형회귀가 나왔습니다. 로지스틱회귀도 같은 순서입니다. y가 0 아니면 1이라서 베르누이를 고릅니다. 와이햇은 y의 기댓값, 베르누이의 기댓값은 파이, 지수족 꼴에서 z 이퀄 로그 오즈를 풀면 파이는 1 나누기 1 더하기 e의 마이너스 z, 3번에 따라 z는 직선입니다. 로지스틱회귀가 나왔습니다. 지난 시간에는 시그모이드를 0과 1 사이의 S 자 함수라서 골랐습니다. 강의 노트의 말대로, 시그모이드가 왜 나왔는지에 대한 답이 여기 있습니다. 고른 게 아니라, y가 0 아니면 1이라서 베르누이를 고른 순간 저절로 나온 함수입니다. 체스로 확인하면, 백이 200점 높을 때 z는 1.03이고 시그모이드를 씌우면 74퍼센트, 지난 시간의 숫자와 같습니다. z를 평균으로 바꾸는 이 함수 g를 정준 반응 함수, canonical response function이라고 부릅니다. 포아송도 같습니다. 와이햇은 평균 람다, z는 log 람다였으니 람다는 e의 z. 앞에서 음수를 피하려고 씌운 e의 z도 고른 게 아니라 포아송에서 나온 함수였습니다.}
\end{frame}

% --- 손실 = -log p
\begin{frame}{손실은 어디서 오나: $-\log p$}
\vfill
\centering
{\LARGE\[ \text{loss} = -\log p(y;z) = \underbrace{-\,z\,y + a(z)}_{\text{\normalsize 분포마다 이 두 항}} \;+\; \text{const} \]}\vskip0.4em
\renewcommand{\arraystretch}{1.45}
{\large\begin{tabular}{l|l}
정규 & $\tfrac12 (y - \hat y)^2$ \;=\; \kb{오차 제곱}\\
베르누이 & $-\big[\,y\log\hat y + (1-y)\log(1-\hat y)\,\big]$ \;=\; \kb{교차 엔트로피}\\
포아송 & $\hat y - y\log\hat y$
\end{tabular}}\\[0.6em]
{\small 최대 가능도: 데이터가 나올 확률을 가장 크게 $=$ $-\log$ 확률의 합을 가장 작게 \quad 범주형 $\to$ $-\log\hat y_{\text{정답}}$}
\vfill
\note{손실은 어디서 왔을까요? GLM에서는 손실을 정하는 것도 요리법의 일부입니다. 학습 데이터가 나올 확률을 가장 크게 하는 w를 찾습니다. 이걸 최대 가능도라고 합니다. 확률을 크게 하는 것은 마이너스 로그 확률을 작게 하는 것과 같으니, 손실은 마이너스 log p입니다. 지수족 꼴에 넣으면 마이너스 z y 더하기 a of z, 그리고 w와 상관없는 상수입니다. 정규분포에서 계산하면 2분의 1 y 빼기 와이햇의 제곱, 1주차의 오차 제곱입니다. 베르누이에서 계산하면 지난 시간의 교차 엔트로피입니다. 그러니까 선형회귀가 오차 제곱을, 로지스틱회귀가 교차 엔트로피를 쓰는 건 따로 고른 게 아니라, 분포를 고르면 같이 정해지는 짝입니다. 포아송이면 와이햇 빼기 y log 와이햇입니다.}
\end{frame}

% --- 기울기 유도
\begin{frame}{기울기가 같은 이유: $a'(z)$ 가 곧 평균}
\vfill
\centering
{\large\begin{align*}
\frac{\partial\,\text{loss}}{\partial z} &= -y + a'(z) && \text{loss} = -z y + a(z) + \text{const}\\[0.3em]
&= -y + \textcolor{ksared}{\hat y} && a'(z) = \text{평균} = \hat y\\[0.3em]
\frac{\partial\,\text{loss}}{\partial w_1} &= (\hat y - y)\,x && z = w_0 + w_1 x
\end{align*}}\vskip0.1em
\renewcommand{\arraystretch}{1.2}
{\large\begin{tabular}{l|c|c|c}
 & $a(z)$ & $a'(z)$ & 평균\\ \hline
베르누이 & $\log(1 + e^{z})$ & $\frac{e^{z}}{1 + e^{z}} = \sigma(z)$ & $\phi$ \;\checkmark\\
정규 & $\tfrac12z^2$ & $z$ & $\mu$ \;\checkmark\\
포아송 & $e^{z}$ & $e^{z}$ & $\lambda$ \;\checkmark
\end{tabular}}
\vfill
\note{이제 네 모델의 기울기가 왜 모두 와이햇 빼기 y였는지 보입니다. 손실 마이너스 z y 더하기 a of z를 z로 미분하면 마이너스 y 더하기 a 프라임 z입니다. 여기서 놀라운 사실이 있습니다. a 프라임 z가 언제나 그 분포의 평균입니다. 아래 표로 확인해 봅시다. 베르누이의 a는 log 1 더하기 e의 z, 미분하면 시그모이드 z, 평균 파이입니다. 정규분포의 a는 2분의 1 z제곱, 미분하면 z, 평균 뮤입니다. 포아송의 a는 e의 z, 미분해도 e의 z, 평균 람다입니다. 평균은 곧 예측 와이햇이니, z로 미분한 값은 와이햇 빼기 y입니다. 마지막으로 z가 w0 더하기 w1 x이니 연쇄 법칙으로 x를 곱하면, w1으로 미분한 값은 와이햇 빼기 y 곱하기 x입니다.}
\end{frame}

% --- 한 번 더 미분: a'' = 분산 → 볼록
\begin{frame}{한 번 더 미분하면: $a''(z)$ 는 분산 $\to$ 손실은 볼록}
\vfill
\centering
{\Large\[ \frac{\partial^2\,\text{loss}}{\partial z^2} = a''(z) = \text{Var}(y) \;\ge\; 0 \]}\vskip0.4em
\renewcommand{\arraystretch}{1.2}
{\large\begin{tabular}{l|c|c}
 & $a''(z)$ & 분산\\ \hline
베르누이 & $\sigma(z)\,\big(1 - \sigma(z)\big)$ & $\phi(1-\phi)$ \;\checkmark\\
정규 & $1$ & $\sigma^2 = 1$ \;\checkmark\\
포아송 & $e^{z}$ & $\lambda$ \;\checkmark
\end{tabular}}\\[0.5em]
{\large 기울기가 늘 0 이상 \;$\to$\; 손실은 \kb{그릇 모양} (볼록) \;$\to$\; 경사하강법이 \kb{하나뿐인 바닥}으로 간다}\\[0.3em]
{\small $z$ 가 $w$ 의 직선이라 $w$ 에 대해서도 볼록 \quad (CS229 강의 4)}
\vfill
\note{한 번 더 미분해 봅시다. 손실을 z로 두 번 미분하면 a 두 번 미분이고, 이건 언제나 y의 분산입니다. 표로 확인하면 베르누이는 시그모이드 곱하기 1 빼기 시그모이드, 파이 곱하기 1 빼기 파이로 베르누이의 분산입니다. 정규는 1, 시그마를 1로 두었으니 분산 1입니다. 포아송은 e의 z, 람다이고 포아송은 분산도 람다입니다. 분산은 음수가 될 수 없으니 두 번 미분한 값이 언제나 0 이상, 손실은 그릇 모양, 볼록 함수입니다. z가 w의 직선이니 w에 대해서도 볼록합니다. 그래서 GLM은 어느 분포를 고르든 바닥이 하나뿐이고, 경사하강법이 엉뚱한 곳에 갇히지 않습니다. 1주차 선형회귀와 지난 시간 로지스틱회귀의 손실이 그릇 모양이었던 이유입니다. 앤드류 응의 강의 4강에서 다루는 지수족의 성질입니다.}
\end{frame}

% --- 한 표
\begin{frame}{정리: 분포만 바꾸면 모델이 나온다}
\vfill
\centering
\renewcommand{\arraystretch}{1.7}
{\large\begin{tabular}{l|l|c|c|l}
$y$ & 분포 & 매개변수 $z$ & 예측 $\hat y = g(z)$ & 모델\\ \hline
실수 & 정규 & $\mu$ & $z$ & 선형회귀\\
0 / 1 & 베르누이 & $\log\frac{\phi}{1-\phi}$ & $\sigma(z)$ & 로지스틱회귀\\
개수 & 포아송 & $\log\lambda$ & $e^{z}$ & 포아송 회귀\\
$K$ 개 중 하나 & 범주형 & (벡터) & $\mathrm{softmax}(z)$ & 소프트맥스 회귀
\end{tabular}}\\[1.0em]
{\large 직선 $z$ 는 그대로 \;·\; 바깥 함수 $g$ 만 바뀐다}\\[0.3em]
{\small $g$ = response function, \;$g^{-1}$ = \kb{link function} (로지스틱: 로그 오즈, 포아송: $\log$)}
\vfill
\note{지금까지를 한 표로 모읍니다. y가 실수면 정규분포, z는 뮤, 예측은 z 그대로, 선형회귀입니다. 0과 1이면 베르누이, z는 로그 오즈, 예측은 시그모이드 z, 로지스틱회귀입니다. 개수면 포아송, z는 log 람다, 예측은 e의 z, 포아송 회귀입니다. 클래스가 여럿이면 범주형, 예측은 소프트맥스, 소프트맥스 회귀입니다. 모든 줄에서 직선 z는 똑같고, 바깥에 씌우는 함수 g만 바뀝니다. 이 g를 반응 함수, 거꾸로 평균을 z로 바꾸는 함수를 링크 함수라고 부릅니다. 로지스틱회귀의 링크 함수는 로그 오즈, 포아송 회귀의 링크 함수는 로그입니다.}
\end{frame}

\section{신경망}
% --- 신경망의 마지막 층
\begin{frame}{신경망으로 이어진다: 마지막 층이 GLM}
\vfill
\centering
\begin{tikzpicture}[node distance=1.1cm, every node/.style={font=\large},
  box/.style={draw, rounded corners=3pt, thick, minimum height=1.0cm, minimum width=1.6cm, align=center}]
\node[box] (x) {입력\\$x$};
\node[box, right=of x, fill=gray!10] (h) {숨은 층들\\특징 $h$};
\node[box, right=of h, draw=ksablue, fill=ksablue!8] (z) {$z = w \cdot h + b$};
\node[box, right=of z, draw=ksablue, fill=ksablue!8] (g) {$\hat y = g(z)$};
\node[box, right=of g] (L) {손실\\$-\log p$};
\draw[-{Stealth},thick] (x) -- (h); \draw[-{Stealth},thick] (h) -- (z); \draw[-{Stealth},thick] (z) -- (g); \draw[-{Stealth},thick] (g) -- (L);
\draw[ksablue, thick, dashed, rounded corners] ([xshift=-6pt,yshift=8pt]z.north west) rectangle ([xshift=6pt,yshift=-8pt]g.south east);
\node[ksablue, font=\large\bfseries, above=10pt of $(z.north)!0.5!(g.north)$] {마지막 층 $=$ GLM};
\draw[-{Stealth}, ultra thick, ksared] ([yshift=-14pt]L.south) to[bend left=18] node[below, font=\large\bfseries] {역전파의 시작: $\hat y - y$} ([yshift=-14pt]z.south);
\end{tikzpicture}\vskip1.0em
{\large 분류 신경망의 끝 $=$ 로지스틱 / 소프트맥스 회귀 \quad 회귀 신경망의 끝 $=$ 선형회귀}\\[0.3em]
{\large 앞 층들은 $x$ 대신 \kb{좋은 특징 $h$} 를 만드는 일을 한다}
\vfill
\note{GLM은 딥러닝으로 그대로 이어집니다. 신경망은 입력 x를 숨은 층들에 통과시켜 특징 h를 만들고, 마지막 층에서 z 이퀄 w 점곱 h 더하기 b를 계산한 뒤 함수 g를 씌웁니다. 이 마지막 층이 바로 오늘 본 GLM입니다. 분류하는 신경망의 끝은 로지스틱회귀나 소프트맥스 회귀이고, 숫자를 맞히는 신경망의 끝은 선형회귀입니다. 앞의 층들이 하는 일은 x 대신 GLM이 잘 쓸 수 있는 좋은 특징 h를 만드는 것입니다. 그리고 학습할 때 손실을 z로 미분한 값은, 오늘 본 대로 와이햇 빼기 y입니다. 역전파는 바로 이 와이햇 빼기 y에서 시작해서 앞 층으로 전해집니다. 7주차 역전파에서 다시 만납니다.}
\end{frame}

% --- 분포 하나에서 출력 함수와 손실이 둘 다
\begin{frame}{분포는 activation 도 loss 도 아니다: 둘 다 분포에서 나온다}
\vfill
\centering
\begin{tikzpicture}[every node/.style={font=\large},
  box/.style={draw, rounded corners=3pt, thick, minimum height=1.0cm, minimum width=3.4cm, align=center}]
\node[box, draw=ksablue, fill=ksablue!8, font=\Large\bfseries] (d) {분포 하나 고르기};
\node[box, below left=1.3cm and 0.6cm of d] (g) {출력 함수 $g(z)$\\{\normalsize 마지막 층의 activation}};
\node[box, below right=1.3cm and 0.6cm of d] (l) {손실 함수\\{\normalsize loss}};
\draw[-{Stealth}, thick] (d.south) -- node[left=6pt, font=\normalsize] {평균을 $z$ 로 풀기} (g.north);
\draw[-{Stealth}, thick] (d.south) -- node[right=6pt, font=\normalsize] {$-\log p$} (l.north);
\node[below=0.25cm of g, font=\normalsize] {예: 베르누이 $\to$ 시그모이드 $\sigma(z)$};
\node[below=0.25cm of l, font=\normalsize] {예: 베르누이 $\to$ 이진 교차 엔트로피};
\end{tikzpicture}\vskip1.2em
{\large 같은 분포에서 나왔으니 \kb{짝}으로 다닌다 \;$\to$\; 미분하면 $\hat y - y$}\\[0.3em]
{\small 숨은 층의 activation (ReLU 등) 은 분포와 상관없다}
\vfill
\note{자주 헷갈리는 질문이 있습니다. 정규분포, 베르누이 같은 분포는 activation 함수일까요, 손실 함수일까요? 둘 다 아닙니다. 분포는 y가 이렇게 나온다는 가정입니다. 그리고 이 가정 하나에서 둘이 다 나옵니다. 분포의 평균을 z로 풀면 출력 함수 g가 나오고, 이게 신경망 마지막 층의 activation입니다. 분포에 마이너스 log를 씌우면 손실 함수가 나옵니다. 베르누이를 고르면 출력은 시그모이드, 손실은 이진 교차 엔트로피입니다. 같은 분포에서 나왔으니 둘은 짝으로 다니고, 그래서 미분하면 와이햇 빼기 y로 깔끔합니다. 단, 숨은 층의 ReLU 같은 activation은 분포와 상관없습니다. 마지막 층만입니다. 다음 쪽에서 네 분포의 짝을 표로 봅니다.}
\end{frame}

% --- 짝 맞추기
\begin{frame}{출력 함수와 손실은 짝으로 쓴다}
\vfill
\centering
\renewcommand{\arraystretch}{1.6}
{\large\begin{tabular}{l|c|l|l}
$y$ & 출력 $g(z)$ & 손실 & PyTorch\\ \hline
실수 & 그대로 & 오차 제곱 & \texttt{MSELoss}\\
0 / 1 & $\sigma$ & 이진 교차 엔트로피 & \texttt{BCEWithLogitsLoss}\\
$K$ 개 중 하나 & softmax & 교차 엔트로피 & \texttt{CrossEntropyLoss}\\
개수 & $e^{z}$ & 포아송 & \texttt{PoissonNLLLoss}
\end{tabular}}\\[1.0em]
{\large 짝이 맞으면 기울기 $= \hat y - y$ \quad \kq{짝이 어긋나면} ($\sigma$ + 오차 제곱): $(\hat y - y)\,\hat y(1-\hat y)$}\\[0.3em]
{\small PyTorch 의 \texttt{WithLogits} · \texttt{CrossEntropyLoss} 는 $g$ 를 손실 안에서 씌운다 $\to$ 마지막 층은 $z$ 만 낸다}
\vfill
\note{그래서 신경망을 만들 때 출력 함수와 손실은 짝으로 고릅니다. y가 실수면 출력은 그대로, 손실은 오차 제곱, PyTorch의 MSELoss입니다. 0과 1이면 시그모이드와 이진 교차 엔트로피, BCEWithLogitsLoss입니다. 여러 클래스면 소프트맥스와 교차 엔트로피, CrossEntropyLoss입니다. 개수면 e의 z와 포아송 손실, PoissonNLLLoss입니다. 짝이 맞으면 기울기가 와이햇 빼기 y로 깔끔합니다. 지난 시간에 본 것처럼 시그모이드에 오차 제곱을 쓰면 시그모이드의 미분 와이햇 곱하기 1 빼기 와이햇이 더 곱해져서, 크게 틀렸는데도 와이햇이 0이나 1 근처면 기울기가 거의 0이 되어 학습이 멈춥니다. 참고로 PyTorch의 WithLogits와 CrossEntropyLoss는 시그모이드나 소프트맥스를 손실 안에서 씌우기 때문에, 신경망의 마지막 층은 z만 내보내면 됩니다.}
\end{frame}

% --- 짝이 어긋나면: 오차 제곱 vs 교차 엔트로피
\begin{frame}{짝이 어긋나면: $\sigma$ 에 오차 제곱 vs 교차 엔트로피}
\vfill
\centering
{\large $z = w \cdot x$, \quad $\hat y = \sigma(z)$, \quad $d\hat y/dz = \hat y(1-\hat y)$}\\[0.2em]
\begin{columns}[T]
\begin{column}{0.47\textwidth}
\centering
{\large\kq{오차 제곱} (Squared Error)}\\[0.3em]
$\displaystyle L = \tfrac12(\hat y - y)^2$\\[0.3em]
$\displaystyle \frac{dL}{d\hat y} = \hat y - y$\\[0.3em]
$\displaystyle \frac{dL}{dw} = (\hat y - y)\cdot\textcolor{ksared}{\hat y(1-\hat y)}\cdot x$\\[0.3em]
{\small $\hat y(1-\hat y)$ 가 \kq{남는다}}
\end{column}
\begin{column}{0.47\textwidth}
\centering
{\large\kb{교차 엔트로피} (Cross Entropy)}\\[0.3em]
$\displaystyle L = -\big[\,y\log\hat y + (1-y)\log(1-\hat y)\,\big]$\\[0.3em]
$\displaystyle \frac{dL}{d\hat y} = \frac{\hat y - y}{\textcolor{ksared}{\hat y(1-\hat y)}}$\\[0.3em]
$\displaystyle \frac{dL}{dw} = (\hat y - y)\cdot x$\\[0.3em]
{\small $\hat y(1-\hat y)$ 가 \kb{지워진다}}
\end{column}
\end{columns}\vskip0.4em
\renewcommand{\arraystretch}{1.15}
{\large\begin{tabular}{l|c}
\multicolumn{2}{l}{정답 $y = 1$ 인데 $\hat y = 0.01$ (크게 틀림) 일 때 기울기 $(dL/dw) / x$}\\ \hline
오차 제곱 & $(0.01 - 1) \times 0.01 \times 0.99 = -0.0098$\\
교차 엔트로피 & $0.01 - 1 = -0.99$ \quad ($100$ 배)
\end{tabular}}
\vfill
\note{짝이 어긋나면 어떻게 되는지 식으로 직접 봅시다. 출력은 시그모이드, z는 w 점곱 x이고, 시그모이드의 미분은 와이햇 곱하기 1 빼기 와이햇입니다. 왼쪽은 오차 제곱입니다. 2분의 1 와이햇 빼기 y의 제곱을 와이햇으로 미분하면 와이햇 빼기 y입니다. 연쇄 법칙으로 시그모이드 미분과 x를 곱하면, 와이햇 빼기 y 곱하기 와이햇 곱하기 1 빼기 와이햇 곱하기 x입니다. 시그모이드 미분이 그대로 남습니다. 오른쪽은 교차 엔트로피, 베르누이의 마이너스 log p입니다. 와이햇으로 미분하면 마이너스 y 나누기 와이햇 더하기 1 빼기 y 나누기 1 빼기 와이햇이고, 통분하면 와이햇 빼기 y 나누기 와이햇 곱하기 1 빼기 와이햇입니다. 분모에 시그모이드 미분이 생겨서, 연쇄 법칙으로 곱하면 정확히 지워지고 와이햇 빼기 y 곱하기 x만 남습니다. 숫자로 보면 차이가 큽니다. 정답이 1인데 와이햇이 0.01로 크게 틀렸을 때, 오차 제곱의 기울기는 마이너스 0.0098, 교차 엔트로피는 마이너스 0.99로 100배 차이입니다. 오차 제곱은 크게 틀렸는데도 거의 움직이지 않습니다. 그래서 출력 함수와 손실을 짝으로 씁니다.}
\end{frame}

\section{연습}
% --- 문제풀이
\begin{frame}{문제풀이}
\vfill
\begin{enumerate}
\item 베르누이에서 $\phi = 0.8$ 이면 $z$ 는? \; ($\log 4 \approx 1.39$)\vskip0.6em
\item 월드컵 모델 $\hat y = e^{\,0.289 + 0.00195\,x}$ ($x$ = Elo 차이). Elo 가 200 앞서면 골은 몇 배? \; ($e^{0.39} \approx 1.48$)\vskip0.6em
\item 클래스 2개 소프트맥스에서 $z_1 = 3$, $z_0 = 1$ 이면 $P(y = 1)$ 은? \; ($\sigma(2) \approx 0.88$)\vskip0.6em
\item 포아송 회귀에서 정답 $y = 5$, 예측 $\hat y = 3$, 입력 $x = 2$. 이 데이터 하나의 $\partial\,\text{loss} / \partial w_1$ 은?\vskip0.6em
\item 일주일 동안 가게에 온 손님 수를 예측하는 신경망. 마지막 층의 출력 함수와 손실은?
\end{enumerate}
\vfill
\note{1번 답. z는 로그 오즈, log 0.8 나누기 0.2, log 4, 약 1.39입니다. 2번 답. 200점 앞서면 z가 0.39 늘고, 골은 e의 0.39, 약 1.48배가 됩니다. 100점마다 1.22배를 두 번 곱한 것과 같습니다. 3번 답. 시그모이드 z1 빼기 z0, 시그모이드 2, 약 0.88입니다. 4번 답. 와이햇 빼기 y는 3 빼기 5, 마이너스 2, 여기에 x 2를 곱해서 마이너스 4입니다. 음수니까 w1을 키워서 예측을 늘리는 쪽으로 갑니다. 포아송이어도 로지스틱회귀와 같은 계산입니다. 5번 답. 손님 수는 개수이니 포아송, 출력 함수는 e의 z, 손실은 포아송 손실, PyTorch의 PoissonNLLLoss입니다.}
\end{frame}

\section{정리}
% --- 정리
\begin{frame}{정리}
\vfill
\centering\large\begin{tabular}{l}
\kb{핵심} \quad 모델 고르기 $=$ $y$ 의 \kb{분포} 가정 $\to$ 씌울 함수 · 손실이 따라 나온다\\[0.5em]
GLM \quad 분포 하나 고르기 + 예측 $=$ 평균 + 매개변수 $=$ 직선 $z$\\[0.5em]
지수족 \quad $p(y;z) = b(y)\exp(z y - a(z))$ \;(정규 · 베르누이 · 포아송 · 범주형)\\[0.5em]
가족 \quad 정규 $\to z$ \;·\; 베르누이 $\to \sigma(z)$ \;·\; 포아송 $\to e^{z}$ \;·\; 범주형 $\to$ softmax\\[0.5em]
학습 \quad 손실 $= -\log p$, 기울기는 언제나 $(\hat y - y)\,x$ \;($a'(z) =$ 평균)\\[0.5em]
볼록 \quad $a''(z) =$ 분산 $\ge 0$ $\to$ 손실은 그릇 모양, 바닥 하나\\[0.5em]
신경망 \quad 마지막 층 $=$ GLM, 역전파는 $\hat y - y$ 에서 시작\\[0.5em]
월드컵 2136줄 \quad Elo 100 앞설 때마다 골 $\times 1.22$, 직선과 달리 음수 없음
\end{tabular}
\vfill
\note{정리합니다. 오늘의 핵심은 모델을 고른다는 것이 y가 어떻게 생겨나는지, 어떤 분포를 따르는지 가정하는 것이고, 그 가정 하나에서 직선에 씌울 함수와 손실이 따라 나온다는 것입니다. GLM은 y의 분포 하나를 고르고, 예측을 그 분포의 평균으로, 분포의 매개변수를 직선 z로 두는 요리법입니다. 정규, 베르누이, 포아송, 범주형 분포는 모두 지수족 꼴로 쓸 수 있습니다. 분포에 따라 예측은 z 그대로, 시그모이드, e의 z, 소프트맥스가 되고, 각각 선형회귀, 로지스틱회귀, 포아송 회귀, 소프트맥스 회귀입니다. 손실은 마이너스 log p이고, a 프라임이 평균이라서 기울기는 언제나 와이햇 빼기 y 곱하기 x이고, a 두 번 미분이 분산이라 손실은 언제나 볼록합니다. 신경망의 마지막 층도 GLM이고, 역전파는 와이햇 빼기 y에서 시작합니다.}
\end{frame}

% --- 참고 문헌
\begin{frame}{참고 문헌}
\vfill
\small
\begin{itemize}
\item Ng, A., \& Ma, T. CS229 Lecture Notes, Chapter 3: Generalized Linear Models. Stanford. \paperlink{https://cs229.stanford.edu/main_notes.pdf}
\item Ng, A. (2018). Lecture 4: Perceptron \& Generalized Linear Model, CS229 Autumn 2018. \paperlink{https://www.youtube.com/watch?v=iZTeva0WSTQ}
\item McCullagh, P., \& Nelder, J. A. (1989). \textit{Generalized Linear Models} (2nd ed.). Chapman \& Hall.
\item Jürisoo, M. International football results from 1872 (국가대표 경기 결과, CC0). Elo 는 World Football Elo Ratings (eloratings.net) 방식으로 직접 계산. \paperlink{https://github.com/martj42/international_results}
\item Fisher, R. A. (1936). The use of multiple measurements in taxonomic problems. \textit{Annals of Eugenics}, 7(2).
\item lichess.org open database, 2013년 1월 (CC0). \paperlink{https://database.lichess.org}
\end{itemize}
\vfill
\note{오늘 내용은 스탠퍼드 CS229 강의 노트의 3부, 일반화 선형 모델을 바탕으로 했습니다. 앤드류 응의 2018년 강의 4강 영상에서 같은 내용을 볼 수 있습니다. 월드컵 데이터는 1872년부터의 국가대표 경기 결과 공개 데이터이고, Elo 점수는 그 경기 기록으로 직접 계산했습니다. 붓꽃은 피셔의 1936년 데이터, 체스는 lichess 공개 데이터베이스입니다.}
\end{frame}
\end{document}
